\subsection{CASE OF THE GROUP OF WEIGHTS; MAXIMAL TERMS}

We retain the notations of the preceding number and let R be a reduced root system in a real vector space V. In the remainder of this section, we take for P the group of weights of R ( \(\S\) 1, no. 9). The group W = W(R) acts on P, hence also on the algebra A{[}P{]}; we have \(w(e^{p}) = e^{w(p)}\) for \(w \in W\) and \(p \in P\) .

Let C be a chamber of R ( \(\S\) 1, no. 5) and let \(B = (\alpha_{i})_{1 \leqslant i \leqslant l}\) be the corresponding basis of R. We provide V (and hence also P) with the order structure defined by C. If \(p, p' \in P\) , \(p \geqslant p'\) if and only if \(p - p'\) is a linear combination of the \(\alpha_{i}\) with positive coefficients.

DEFINITION 1. Let \(x = \sum_{p \in \mathrm{P}} x_p e^p\) be an element of A{[}P{]}. The set S of \(p \in \mathrm{P}\) such that \(x_p \neq 0\) is called the support of \(x\) and the set X of maximal elements of S is called the maximal support of \(x\) . A term \(x_p e^p\) with \(p \in X\) is then called a maximal term of \(x\) .

Lemma 2. Let \(x \in \mathrm{A}[\mathrm{P}]\) and let \((x_{p}e^{p})_{p \in \mathrm{X}}\) be the family of maximal terms of \(x\) . Let \(q \in \mathrm{P}\) and let \(y \in \mathrm{A}[\mathrm{P}]\) be such that \(e^{q}\) is the unique maximal term of \(y\) . Then, the family of maximal terms of \(xy\) is \((x_{p}e^{p + q})_{p \in \mathrm{X}}\) .

Put \(x = \sum_{p} x_{p} e^{p}, y = \sum_{r} y_{r} e^{r}\) and \(xy = \sum_{t} z_{t} e^{t}\) . Then \(r \leqslant q\) for all \(r \in \mathbb{P}\) such that \(y_{r} \neq 0\) and \(z_{t} = \sum_{p + r = t} x_{p} y_{r}\) .

If \(t = p + q = p' + r\) with \(p \in \mathrm{X}\) and \(x_{p'}y_r \neq 0\) , then \(r \leqslant q\) , hence \(p' \geqslant p\) and consequently \(p' = p\) . Thus \(z_{p + q} = x_p y_q = x_p \neq 0\) . This shows that \(\mathrm{X} + q\) is contained in the support of the product \(xy\) .

On the other hand, if \(t = p' + r\) with \(x_{p'} y_r \neq 0\) , there exists \(p \in X\) such that \(p' \leqslant p\) and we have \(t \leqslant p + q\) . The maximal support of xy is therefore contained in \(X + q\) . Since no two elements of \(X + q\) are comparable, it follows that \(X + q\) is exactly the maximal support of xy and we have seen above that \(z_{p+q} = x_p\) for \(p \in X\) , which completes the proof of the lemma.

Remark. Since \(x \neq 0\) means that the maximal support of x is non-empty, Lemma 2 shows that \(x \neq 0\) implies \(xy \neq 0\) whenever y admits a unique maximal term of the form \(e^{q}\) .

